\section{Introduction}
\label{sec:intro}
Speculative decoding proposes several tokens with a draft model and verifies them with a target, reducing the number of sequential target invocations when proposals are accepted~\cite{leviathan2023}. The execution change is not merely a reduction in calls. Verification processes several query positions together, while ordinary greedy decoding processes one continuation position at a time. In a finite-precision implementation, two mathematically equivalent organizations of that work can produce different scores and, at a close decision boundary, different tokens. A correct acceptance protocol therefore does not by itself establish reproduction of a particular deployed greedy reference.

This distinction matters when a runtime optimization must preserve an existing regression baseline. We study that explicit requirement: the target weights, input formatting, sampler, stopping rule and ordinary execution recipe define a registered reference. A candidate is eligible for a reference-preserving latency comparison only when its complete native output satisfies the same requirement. We do not equate reference preservation with answer correctness, statistical sampling equivalence, or equality of every internal floating-point value. An implementation can be repeatable yet consistently disagree with this reference.

A direct response is to serialize the target computation. That removes opportunities for batching well beyond the operations that need numerical control. Our central observation is that \emph{the organization of attention queries can be separated from the batching of the surrounding target network}. Attention can retain one query column per batch instance while projection, output projection and feed-forward operators continue to process several tokens together. This is a concrete graph-design choice: it changes the attention representation without introducing a new draft algorithm or a new arithmetic kernel.

We implement this choice as a \emph{phase-scoped query layout}. After the runtime's existing permutations, the query view changes from $(d,q,h, 1)$ to $(d, 1,h,q)$. A corresponding mask view preserves each query's admissible keys, and shared K/V remain associated with the same logical stream. The method also includes two runtime obligations. A role/phase guard activates the layout in the intended target phases even when observation is disabled; a graph-reuse key prevents a cached unconstrained topology from bypassing that guard. These obligations are part of the method, not incidental instrumentation.

The empirical argument separates three questions. First, does the transformation have the intended numerical behavior on identical attention inputs? Second, does the complete decoder preserve its registered output? Third, is this implementation cheaper than another implementation that passes the same output test? The first two cannot substitute for each other: identical sampled operators do not establish a whole-model theorem, and matching tokens can conceal numerical differences. The third requires more than comparing a candidate with an unoptimized ordinary decoder. We compare query-batch attention with explicit per-query attention and target-global serialization, then separately measure whether adding a draft is worthwhile on the repaired path.

The CUDA results support all three questions within the tested scope. In 24 captured contexts, query-batch and per-query attention agree bitwise at QK, softmax and AV, whereas ordinary multi-query execution differs. A larger confirmation campaign checks Qwen3, SmolLM2 and Llama under a frozen target-phase configuration. A subsequent implementation campaign executes 2,016 complete requests across seven arms, with all arms meeting the common output requirement. Query-batch is cheaper in aggregate than per-query attention and global serialization in each family, although some per-task comparisons regress and the per-query advantage is small. Drafting remains unprofitable for SmolLM2. These are distinct findings: a less expensive repair need not make speculation profitable.

A separate Vulkan investigation tests the limit of an attention-centered explanation. Changing prefill partitioning at fixed context settings reproduces two historical output divergences. First-layer projection differences appear before later attention, and a validated boundary-state reconstruction shows that the affected suffix follows the retained KV source. This experiment is not a transfer of the CUDA query-batch patch. Instead, it explains why a layout method must state its admitted runtime conditions and validate state formation, rather than assume that attention is the only numerically sensitive operation on every backend.

The paper makes three contributions. \textbf{(1) A specified runtime transformation:} query and mask address mappings, K/V and grouped-query semantics, output ordering, phase activation and cache-reuse conditions define an implementable attention layout. \textbf{(2) Layered numerical and output evidence:} fixed-input operator replay and three-family complete-output qualification establish different, explicitly bounded properties. \textbf{(3) Same-reference implementation and state-boundary evaluation:} fair per-query/query-batch/serialized comparisons establish repair cost separately from draft value, while Vulkan partition and KV interventions delimit portability. We make no claim of a new sampling rule, a universal minimal constraint, or superiority to an external determinism system under a common workload.
